\documentclass[12pt]{article}
% \usepackage[utf8]{inputenc}
\usepackage{graphicx}
\usepackage{listings}

\setlength{\parindent}{0em}

\def\code#1{\texttt{#1}}

% FILL IN TITLE AND AUTHOR
\title{
{CS3102 P2: Practical Report}\\
{\large Reliable Data Transfer Using UDP}\\
{\includegraphics[width=80mm]{university-logo.png}}
}
\author{190010906}
\date{01 April 2022}

\begin{document}

\maketitle

\newpage

\section{Introduction}

This report cover the design and implementation of a connetion-oritented, reliable, unicast, transport protocol,built on top of UDP.

The protocol in question is called RDT - Reliable Data Transport

\section{Design}

This section will describe the design of the RDT protocol and the considerations that informed this design.

\subsection{Packet Structure}

RDT packets are constructed with the following structure:

\begin{figure}[h]
\begin{verbatim}
 0                   1                   2                   3  
 0 1 2 3 4 5 6 7 8 9 0 1 2 3 4 5 6 7 8 9 0 1 2 3 4 5 6 7 8 9 0 1
+-+-+-+-+-+-+-+-+-+-+-+-+-+-+-+-+-+-+-+-+-+-+-+-+-+-+-+-+-+-+-+-+
|            Sequence           |              Type             |
+-+-+-+-+-+-+-+-+-+-+-+-+-+-+-+-+-+-+-+-+-+-+-+-+-+-+-+-+-+-+-+-+
|              Size             |            Checksum           |
+-+-+-+-+-+-+-+-+-+-+-+-+-+-+-+-+-+-+-+-+-+-+-+-+-+-+-+-+-+-+-+-+
\end{verbatim}
\caption{RDT Header}
\end{figure}

The \texttt{type} field contains one following values, denoting the type of the packet:

\begin{itemize}
    \item 0 - SYN.
    \item 1 - SYN ACK
    \item 2 - DATA
    \item 3 - DATA ACK 
    \item 4 - FIN 
    \item 5 - FIN ACK
\end{itemize}

This approach was chosen over a flag-based approach, as it makes it easier to check packet type and there were only a small number of types to define given the simple nature of the protocol.

\subsection{Finite State Machine}

\section{Testing}

This section will detail how RDT was tested to validate correct operation.

\subsection{Methodology}

To test the ability of RDT to deliver packets in a reliable and ordered manner, two test programs were created. The latter was run on \textbf{pc} and the former on \textbf{pc}. Slurpe was place in the middle. A file was transmitted from A to B. Decoded with SHA.


\end{document}